\section{Switch Statements}

Switch statements follow specific formatting rules to ensure consistency and readability throughout the STAR firmware codebase.

\subsection{General Rules}

\begin{itemize}
    \item \textbf{Same Line for Switch and Opening Brace}: The opening brace \texttt{\{} for a \texttt{switch} statement should be on the same line as the \texttt{switch} keyword.
    \item \textbf{Braces for Case Blocks}: Each \texttt{case} should have its own block of code enclosed in braces \texttt{\{\}}.
    \item \textbf{Break Statements}: Every case must end with \texttt{break}, \texttt{return}, or a fall-through comment.
    \item \textbf{Default Case}: Always include a \texttt{default} case, even if it only logs a warning.
\end{itemize}

\subsection{Basic Switch Structure}

\begin{lstlisting}[language=C]
/* CORRECT - from star_bus_manager.c */
switch (config->type) {
    case k_star_bus_type_i2c: {
        ret = internal_register_bus_pin(pin_iface,
                                        config->proto.i2c.config.sda_io_num,
                                        config->name,
                                        "I2C SDA",
                                        true);
        if (ret != ESP_OK && first_error == ESP_OK) {
            first_error = ret;
        }

        ret = internal_register_bus_pin(pin_iface,
                                        config->proto.i2c.config.scl_io_num,
                                        config->name,
                                        "I2C SCL",
                                        true);
        if (ret != ESP_OK && first_error == ESP_OK) {
            first_error = ret;
        }
        break;
    }

    case k_star_bus_type_spi: {
        ret = internal_register_bus_pin(pin_iface,
                                        config->proto.spi.copi_io_num,
                                        config->name,
                                        "SPI COPI",
                                        true);
        if (ret != ESP_OK && first_error == ESP_OK) {
            first_error = ret;
        }
        /* ... more pin registrations ... */
        break;
    }

    default: {
        ESP_LOGD(s_TAG, "No pin registration for bus type: %d", config->type);
        break;
    }
}
\end{lstlisting}

\subsection{Case Block Braces}

Each case must have its own braces. This prevents variable scope issues and improves readability.

\begin{lstlisting}[language=C]
/* CORRECT - Each case has braces */
switch (bus_type) {
    case k_star_bus_type_i2c: {
        i2c_port_t port = config->proto.i2c.port;  /* Local variable */
        uint8_t address = config->proto.i2c.address;
        ESP_LOGI(s_TAG, "I2C bus on port %d, address 0x%02X", port, address);
        break;
    }
    case k_star_bus_type_spi: {
        spi_host_device_t host = config->proto.spi.host;  /* Different scope */
        ESP_LOGI(s_TAG, "SPI bus on host %d", host);
        break;
    }
    case k_star_bus_type_gpio: {
        uint8_t pin_count = config->proto.gpio.pin_count;
        ESP_LOGI(s_TAG, "GPIO bus with %d pins", pin_count);
        break;
    }
    default: {
        ESP_LOGW(s_TAG, "Unknown bus type: %d", bus_type);
        break;
    }
}

/* WRONG - Missing braces */
switch (bus_type) {
    case k_star_bus_type_i2c:
        handle_i2c();
        break;
    case k_star_bus_type_spi:
        handle_spi();
        break;
}
\end{lstlisting}

\subsection{Break Statement Requirements}

Every case must have a clear exit: \texttt{break}, \texttt{return}, or explicit fall-through.

\begin{lstlisting}[language=C]
/* CORRECT - Using break */
switch (config->type) {
    case k_star_bus_type_i2c: {
        process_i2c(config);
        break;
    }
    case k_star_bus_type_spi: {
        process_spi(config);
        break;
    }
    default: {
        break;
    }
}

/* CORRECT - Using return */
const char* star_bus_type_to_string(star_bus_type_t type)
{
    switch (type) {
        case k_star_bus_type_none: {
            return "NONE";
        }
        case k_star_bus_type_i2c: {
            return "I2C";
        }
        case k_star_bus_type_spi: {
            return "SPI";
        }
        case k_star_bus_type_gpio: {
            return "GPIO";
        }
        case k_star_bus_type_adc: {
            return "ADC";
        }
        case k_star_bus_type_count: {
            return "COUNT";
        }
        default: {
            return "UNKNOWN";
        }
    }
}

/* CORRECT - Explicit fall-through (rare, must be commented) */
switch (priority) {
    case PRIORITY_CRITICAL: {
        log_critical(message);
        /* FALLTHROUGH */
    }
    case PRIORITY_HIGH: {
        notify_admin(message);
        /* FALLTHROUGH */
    }
    case PRIORITY_NORMAL: {
        log_message(message);
        break;
    }
    default: {
        break;
    }
}
\end{lstlisting}

\subsection{Default Case Handling}

Always include a default case. Use it for error handling or logging unexpected values.

\begin{lstlisting}[language=C]
/* CORRECT - Default case logs warning */
switch (curr->type) {
    case k_star_bus_type_i2c: {
        reg_result = register_pin_if_valid(curr->proto.i2c.config.sda_io_num,
                                           bus_name, "I2C SDA", reg_func,
                                           user_ctx, true);
        /* ... */
        break;
    }
    case k_star_bus_type_spi: {
        reg_result = register_pin_if_valid(curr->proto.spi.copi_io_num,
                                           bus_name, "SPI COPI", reg_func,
                                           user_ctx, true);
        /* ... */
        break;
    }
    default: {
        ESP_LOGW(manager->tag, "Skipping pin registration for "
                               "unsupported bus type: %d", curr->type);
        break;
    }
}

/* CORRECT - Default case returns error for invalid input */
esp_err_t process_bus_type(star_bus_type_t type)
{
    switch (type) {
        case k_star_bus_type_i2c: {
            return process_i2c();
        }
        case k_star_bus_type_spi: {
            return process_spi();
        }
        default: {
            ESP_LOGE(s_TAG, "Invalid bus type: %d", type);
            return ESP_ERR_INVALID_ARG;
        }
    }
}
\end{lstlisting}

\subsection{Switch on Enum Types}

When switching on an enum, consider handling all values explicitly.

\begin{lstlisting}[language=C]
/* CORRECT - All enum values handled */
switch (config->type) {
    case k_star_bus_type_none: {
        ESP_LOGW(s_TAG, "Bus type is NONE - skipping");
        break;
    }
    case k_star_bus_type_i2c: {
        /* Initialize I2C */
        break;
    }
    case k_star_bus_type_spi: {
        /* Initialize SPI */
        break;
    }
    case k_star_bus_type_gpio: {
        /* Initialize GPIO */
        break;
    }
    case k_star_bus_type_adc: {
        /* Initialize ADC */
        break;
    }
    case k_star_bus_type_count: {
        /* Sentinel value - should never be used */
        ESP_LOGE(s_TAG, "Invalid bus type: count sentinel");
        return ESP_ERR_INVALID_ARG;
    }
    /* No default needed when all enum values are handled */
    /* Compiler warns if new enum values are added */
}
\end{lstlisting}

\subsection{Indentation}

Case labels are at the same indentation level as the switch. Case body is indented one level.

\begin{lstlisting}[language=C]
/* CORRECT indentation */
switch (state) {
    case STATE_IDLE: {
        /* Case body indented */
        if (condition) {
            /* Nested code further indented */
        }
        break;
    }
    case STATE_RUNNING: {
        process_running();
        break;
    }
    default: {
        break;
    }
}

/* WRONG - Case labels indented */
switch (state) {
        case STATE_IDLE: {
            break;
        }
}
\end{lstlisting}

\subsection{Summary}

\begin{itemize}
    \item Opening brace on same line as \texttt{switch}
    \item Each \texttt{case} has its own braces \texttt{\{...\}}
    \item Every case ends with \texttt{break}, \texttt{return}, or fall-through comment
    \item Always include \texttt{default} case
    \item Case labels at switch indentation level
    \item Case body indented one level inside braces
\end{itemize}

